% Volume VIII: Computational Neuroscience Mathematics
% Continuity note for Chapter 44.
% Previous dependencies: Chapter 43's membrane equations and conductance models; Chapter 25's ODE, fixed point, stability, and phase-plane language.
% New durable concepts: threshold-reset models, leaky integrate-and-fire dynamics, interspike intervals, firing-rate curves, phase planes, nullclines, Jacobian classification, adaptation variables, bursting by time-scale separation, and population rate equations.
% Forward hooks: Chapter 45 uses spike times to drive synapses and plasticity; Chapter 46 treats spike times statistically; Chapter 49 reuses rate models in recurrent circuits, attractors, oscillations, and neural fields.
% Notation commitments: point-neuron voltage is V(t); threshold, reset, and refractory values are V_th, V_reset, and tau_ref; firing rate is r; adaptation variables are w or a; population rates are r_i(t) or vector r(t).
\section{Chapter 44: Point Neuron Models and Spiking Dynamics}
\label{ch:point-neuron-models-spiking-dynamics}

Chapter 43 modeled the membrane with physical conductances. For many circuit and learning questions, however, the exact action-potential waveform is not the main object. We often need a model that predicts when spikes occur, how firing rate depends on input, how voltage and recovery variables shape excitability, and how large populations can be reduced to rates.

\paragraph{Chapter problem.}
How can we replace detailed channel biophysics by mathematically controlled point-neuron and rate models while preserving the spike timing, stability, adaptation, and population dynamics needed for computation?

\paragraph{Section map.}
Section 44.1 introduces integrate-and-fire models with threshold, reset, refractory period, and firing-rate curves. Section 44.2 develops phase-plane analysis for two-variable neuron models. Section 44.3 adds slow adaptation variables and explains bursting through time-scale separation. Section 44.4 derives population rate models as coarse descriptions of many spiking neurons.

\subsection{44.1 Integrate-and-fire models}

\paragraph{Guiding question.}
How can a neuron emit discrete spikes if its subthreshold voltage obeys a continuous ODE?

\paragraph{Beginner-facing intuition.}
The leaky integrate-and-fire model keeps the passive membrane equation between spikes and replaces the complicated spike-generating conductances by an event rule. Voltage integrates input through a leaky capacitor. When voltage reaches a threshold, the model records a spike, resets the voltage, and optionally waits through a refractory period.

This hybrid structure is the point. The subthreshold dynamics are continuous and differentiable. The spike event is discrete. The model therefore lives between ODEs and event processes, which makes it useful for networks where spike times matter but channel details do not.

\paragraph{Formal definitions and notation.}
Let \(V(t)\in\mathbb{R}\) be voltage, \(E_L\) the leak reversal potential, \(R_m\) input resistance, \(\tau_m\) membrane time constant, and \(I(t)\) input current. The leaky integrate-and-fire (LIF) subthreshold equation is
\[
\tau_m\frac{dV}{dt}=-(V-E_L)+R_m I(t),
\qquad V(t)<V_{\mathrm{th}}.
\]
The event rule is
\[
V(t^-)=V_{\mathrm{th}}
\quad\Longrightarrow\quad
\text{record a spike at }t,\quad V(t^+)=V_{\mathrm{reset}},
\]
followed by an optional refractory interval \(\tau_{\mathrm{ref}}\) during which no spike is emitted.

The spike train produced by the model is the ordered set of spike times
\[
0<t_1<t_2<\cdots,
\]
or, in generalized-function notation introduced later in Chapter 46,
\[
s(t)=\sum_k\delta(t-t_k).
\]

\paragraph{Core derivation: firing rate under constant current.}
Let \(I(t)=I_0\) and define
\[
V_\infty=E_L+R_m I_0.
\]
Starting from \(V(0)=V_{\mathrm{reset}}\), the subthreshold solution is
\[
V(t)=V_\infty+\bigl(V_{\mathrm{reset}}-V_\infty\bigr)e^{-t/\tau_m}.
\]
If \(V_\infty\leq V_{\mathrm{th}}\), threshold is never reached and the firing rate is \(0\). If \(V_\infty>V_{\mathrm{th}}\), set \(V(T)=V_{\mathrm{th}}\):
\[
V_{\mathrm{th}}=V_\infty+\bigl(V_{\mathrm{reset}}-V_\infty\bigr)e^{-T/\tau_m}.
\]
Solving gives
\[
T=\tau_m\log\!\left(
\frac{V_\infty-V_{\mathrm{reset}}}{V_\infty-V_{\mathrm{th}}}
\right).
\]
Including refractoriness,
\[
\boxed{
r(I_0)=
\left[
\tau_{\mathrm{ref}}
+\tau_m\log\!\left(
\frac{V_\infty-V_{\mathrm{reset}}}{V_\infty-V_{\mathrm{th}}}
\right)
\right]^{-1}
}
\]
for \(V_\infty>V_{\mathrm{th}}\), and \(r(I_0)=0\) otherwise. This is the LIF \(f\)-\(I\) curve.

\paragraph{Concrete example.}
Let \(E_L=-65\) mV, \(V_{\mathrm{reset}}=-70\) mV, \(V_{\mathrm{th}}=-50\) mV, \(\tau_m=20\) ms, \(\tau_{\mathrm{ref}}=2\) ms, and \(R_m I_0=25\) mV. Then \(V_\infty=-40\) mV, so threshold is reachable. The integration time is
\[
T=20\log\!\left(\frac{-40-(-70)}{-40-(-50)}\right)
=20\log(3)\approx 22.0\ \mathrm{ms}.
\]
The total interspike interval is \(24.0\) ms, so
\[
r\approx 41.7\ \mathrm{Hz}.
\]
If \(R_m I_0=10\) mV, then \(V_\infty=-55\) mV, below threshold, and the model emits no spikes.

\paragraph{Computational neuroscience example.}
Large spiking-network simulations often use LIF neurons because they expose the computational role of synaptic input, recurrent connectivity, and spike timing without paying the cost of Hodgkin-Huxley gates. Balanced excitatory-inhibitory networks, winner-take-all circuits, and sensory encoding models can all be studied with threshold-reset units.

\paragraph{AI and machine-learning example.}
Spiking neural networks use integrate-and-fire units to process event streams and to target neuromorphic hardware. The threshold-reset operation is not differentiable in the ordinary sense, so training often uses surrogate gradients: the forward pass emits spikes, while the backward pass replaces the hard threshold derivative with a smooth approximation. This is a modeling choice, not a biological law.

\paragraph{Common misconceptions and failure modes.}
First, the LIF spike has no waveform; it records an event time. Do not interpret the reset jump as a biophysical voltage trajectory. Second, threshold is a rule in this model, not an emergent conductance instability. Third, firing rate formulas assume constant input and regular spiking; noisy input produces variable interspike intervals. Fourth, the refractory period sets an upper rate but does not by itself model all post-spike channel recovery.

\paragraph{Exercises.}
\begin{enumerate}[label=\textbf{44.1.\arabic*.}, leftmargin=*]
\item \textbf{Mechanical.} Identify \(E_L\), \(V_{\mathrm{th}}\), \(V_{\mathrm{reset}}\), \(\tau_m\), and \(R_m I(t)\) in the LIF equation.
\item \textbf{Proof or derivation.} Derive the LIF interspike interval under constant current from the exponential subthreshold solution.
\item \textbf{Numerical/computational.} With \(E_L=-65\) mV, \(V_{\mathrm{reset}}=-65\) mV, \(V_{\mathrm{th}}=-50\) mV, \(\tau_m=10\) ms, \(\tau_{\mathrm{ref}}=1\) ms, and \(R_m I_0=25\) mV, compute the firing rate.
\item \textbf{Interpretation.} Why does increasing leak conductance usually shorten the integration window for coincident synaptic inputs?
\item \textbf{Misconception/debugging.} A student plots \(V(t)\) from an LIF model and expects to see action-potential peaks. Explain what the model records instead.
\end{enumerate}

\subsection{44.2 Phase-plane analysis of neurons}

\paragraph{Guiding question.}
How do voltage and a recovery variable interact geometrically to produce rest, threshold-like behavior, and repetitive spiking?

\paragraph{Beginner-facing intuition.}
One voltage variable can relax or cross a threshold, but many spike mechanisms require at least two interacting variables: a fast voltage-like variable and a slower recovery variable. The phase plane draws the state \((V,w)\) as a point. Arrows show how the point moves. Nullclines show where one component temporarily stops changing. Fixed points, separatrices, and limit cycles become visible geometry.

\paragraph{Formal definitions and notation.}
A two-variable neuron model has the form
\[
\frac{dV}{dt}=f(V,w;I),\qquad
\frac{dw}{dt}=g(V,w;I),
\]
where \(V\) is voltage or voltage-like activity, \(w\) is a recovery variable, and \(I\) is an input parameter. The phase plane is the \((V,w)\)-plane.

The \(V\)-nullcline is
\[
f(V,w;I)=0,
\]
and the \(w\)-nullcline is
\[
g(V,w;I)=0.
\]
A fixed point \((V^*,w^*)\) satisfies both equations. The Jacobian at a fixed point is
\[
J=
\begin{bmatrix}
\partial f/\partial V & \partial f/\partial w\\
\partial g/\partial V & \partial g/\partial w
\end{bmatrix}_{(V^*,w^*)}.
\]
Its eigenvalues determine local stability. For a \(2\times2\) matrix, the trace \(\tau=\operatorname{tr}(J)\) and determinant \(\Delta=\det(J)\) give
\[
\lambda_{1,2}=\frac{\tau\pm\sqrt{\tau^2-4\Delta}}{2}.
\]

\paragraph{Core mechanism: nullclines and linearization.}
Nullclines organize the flow. On the \(V\)-nullcline, the vector field is vertical because \(dV/dt=0\). On the \(w\)-nullcline, the vector field is horizontal because \(dw/dt=0\). Their intersections are resting or unstable states.

Near a fixed point, write \(z=(V,w)^\top\) and \(z^*=(V^*,w^*)^\top\). Taylor approximation gives
\[
\frac{d}{dt}(z-z^*)\approx J(z-z^*).
\]
If all eigenvalues of \(J\) have negative real parts, small perturbations decay. If some eigenvalue has positive real part, perturbations grow in that direction. In neuron models, an unstable fixed point or saddle can help form a threshold separatrix, while a stable limit cycle can represent repetitive spiking.

\paragraph{Concrete example.}
Suppose a reduced neuron model has a fixed point whose Jacobian is
\[
J=
\begin{bmatrix}
-0.5 & -1.0\\
0.2 & -0.3
\end{bmatrix}.
\]
Then \(\tau=-0.8\) and
\[
\Delta=(-0.5)(-0.3)-(-1)(0.2)=0.35.
\]
The discriminant is
\[
\tau^2-4\Delta=0.64-1.40=-0.76<0.
\]
The eigenvalues are complex with real part \(\tau/2=-0.4\), so the fixed point is a stable spiral. A voltage perturbation returns to rest with damped oscillations. In a neuron, this corresponds to a resonant subthreshold response rather than monotone relaxation.

\paragraph{Computational neuroscience example.}
The FitzHugh-Nagumo model is a common two-variable caricature:
\[
\dot V=V-\frac{V^3}{3}-w+I,\qquad
\dot w=\varepsilon(V+a-bw).
\]
The cubic \(V\)-nullcline and approximately linear \(w\)-nullcline explain excitability: a stimulus can push the state past a knee of the cubic, after which the trajectory makes a large excursion before returning. This gives a geometric explanation of spike initiation without simulating every ion channel.

\paragraph{AI and machine-learning example.}
Continuous-time recurrent neural networks and neural ODEs also have phase portraits. A two-dimensional hidden state can have stable points, spirals, or oscillations depending on its Jacobian and nonlinearities. The same phase-plane tools diagnose whether a learned recurrent module stores a memory, damps perturbations, oscillates, or becomes unstable.

\paragraph{Common misconceptions and failure modes.}
First, a nullcline is not a trajectory. It is a curve where one velocity component is zero. Second, threshold in a phase plane is often a curve or separatrix, not a single voltage. Third, local linearization describes behavior near a fixed point; it cannot by itself prove the existence of a large-amplitude spike or limit cycle. Fourth, complex eigenvalues mean rotation in the linearized system, not necessarily sustained oscillation.

\paragraph{Exercises.}
\begin{enumerate}[label=\textbf{44.2.\arabic*.}, leftmargin=*]
\item \textbf{Mechanical.} For \(\dot V=f(V,w)\), \(\dot w=g(V,w)\), define the \(V\)-nullcline, \(w\)-nullcline, and fixed point.
\item \textbf{Proof or derivation.} Derive the characteristic equation \(\lambda^2-\tau\lambda+\Delta=0\) for a \(2\times2\) Jacobian with trace \(\tau\) and determinant \(\Delta\).
\item \textbf{Numerical/computational.} For \(J=\begin{bmatrix}0.2&-1\\0.5&-0.4\end{bmatrix}\), compute trace, determinant, discriminant, and classify local stability.
\item \textbf{Interpretation.} Why does a slow recovery variable help a trajectory return after the voltage upstroke?
\item \textbf{Misconception/debugging.} A student says, ``If a point lies on the \(V\)-nullcline, the system is stopped.'' Explain why the \(w\)-component can still move.
\end{enumerate}

\subsection{44.3 Spike-frequency adaptation and bursting}

\paragraph{Guiding question.}
How can a neuron change its firing rate during sustained input or produce clusters of spikes separated by silence?

\paragraph{Beginner-facing intuition.}
Real neurons do not always fire like a metronome under constant input. Many slow down after the first few spikes. Some fire bursts: several spikes in a row, then a quiet interval, then another burst. A common mathematical explanation is a slow variable that accumulates during spiking and relaxes during silence. This slow variable changes the effective input or threshold, creating time-dependent excitability.

\paragraph{Formal definitions and notation.}
An adaptive integrate-and-fire model adds a slow adaptation current \(w(t)\):
\[
C\frac{dV}{dt}=-g_L(V-E_L)-w+I(t),
\]
\[
\tau_w\frac{dw}{dt}=a(V-E_L)-w,
\]
with a spike-triggered jump
\[
V(t^-)=V_{\mathrm{th}}
\quad\Longrightarrow\quad
\text{spike},\quad V(t^+)=V_{\mathrm{reset}},\quad w(t^+)=w(t^-)+b.
\]
Here \(a\geq0\) controls subthreshold adaptation, \(b\geq0\) controls spike-triggered adaptation, and \(\tau_w\) is usually larger than \(\tau_m\).

Bursting requires at least one slow process that moves the fast spiking subsystem between active and silent regimes. A schematic fast-slow form is
\[
\dot x=f(x,y;\mu),\qquad
\dot y=\varepsilon g(x,y;\mu),\qquad 0<\varepsilon\ll1,
\]
where \(x\) is fast and \(y\) is slow.

\paragraph{Core mechanism: adaptation as negative feedback.}
During repetitive firing, each spike increases \(w\) by \(b\). The voltage equation contains \(-w\), so increasing \(w\) reduces net depolarizing drive. Between spikes and after input ends, \(w\) relaxes toward its voltage-dependent baseline. Therefore the effective drive is high at stimulus onset and lower after several spikes.

If spike \(k\) occurs at time \(t_k\), and if we ignore the subthreshold \(a(V-E_L)\) term for a simple approximation, then between spikes
\[
w(t)=w(t_k^+)e^{-(t-t_k)/\tau_w}.
\]
At each spike,
\[
w(t_k^+)=w(t_k^-)+b.
\]
Thus a high firing rate makes \(w\) accumulate. The accumulation raises the interspike interval, giving spike-frequency adaptation.

\paragraph{Concrete example.}
Let \(b=40\) pA and \(\tau_w=200\) ms. Suppose a neuron fires every \(20\) ms, and ignore subthreshold drive into \(w\). The decay factor between spikes is \(e^{-20/200}=e^{-0.1}\approx0.905\). If \(w\) has reached a steady cycle just before a spike, then
\[
w^- = (w^-+b)e^{-0.1}.
\]
Solving,
\[
w^-=\frac{b e^{-0.1}}{1-e^{-0.1}}\approx
\frac{40(0.905)}{0.095}\approx381\ \mathrm{pA}.
\]
The after-spike value is about \(421\) pA. This large opposing current can substantially reduce the firing rate compared with the first spike.

\paragraph{Computational neuroscience example.}
Spike-frequency adaptation can arise from slow potassium currents, calcium-activated potassium currents, sodium-channel inactivation, or synaptic mechanisms. It affects coding: an adapting neuron responds strongly to changes or stimulus onset but less strongly to sustained input. Bursting can arise from calcium currents interacting with slower potassium or calcium-dependent processes, producing alternating active and silent phases.

\paragraph{AI and machine-learning example.}
Adaptation variables appear in spiking neural networks as dynamic thresholds or fatigue terms. They can prevent runaway firing, encode recent activity history, and improve temporal feature extraction. In recurrent machine-learning models, similar slow variables act as memory or gating states. The precise biological interpretation should not be assumed unless the variable is tied to a measured current or mechanism.

\paragraph{Common misconceptions and failure modes.}
First, adaptation is not the same as refractoriness. Refractoriness acts immediately after a spike; adaptation can accumulate over many spikes and decay slowly. Second, bursting is not just ``high firing rate.'' It is a structured alternation of active and silent phases. Third, adding a slow variable can change stability and bifurcations, not merely shift the firing-rate curve. Fourth, many biological mechanisms can produce similar adaptation traces, so model identifiability is limited without additional data.

\paragraph{Exercises.}
\begin{enumerate}[label=\textbf{44.3.\arabic*.}, leftmargin=*]
\item \textbf{Mechanical.} In the adaptive integrate-and-fire equations, identify which terms produce leak, input, subthreshold adaptation, and spike-triggered adaptation.
\item \textbf{Proof or derivation.} Derive the steady pre-spike adaptation value \(w^-=(w^-+b)e^{-\Delta/\tau_w}\) for periodic spiking with interval \(\Delta\), and solve for \(w^-\).
\item \textbf{Numerical/computational.} For \(b=20\) pA, \(\tau_w=100\) ms, and \(\Delta=10\) ms, compute the steady pre-spike \(w^-\) in the simplified jump-decay model.
\item \textbf{Interpretation.} Why can adaptation make a neuron sensitive to stimulus changes rather than sustained stimulus level?
\item \textbf{Misconception/debugging.} A model fires one spike at stimulus onset and then stops. Give two adaptation-related parameter changes that could restore sustained firing.
\end{enumerate}

\subsection{44.4 Population rate models}

\paragraph{Guiding question.}
When many neurons spike, when is it useful to model their collective activity by firing rates rather than individual spike times?

\paragraph{Beginner-facing intuition.}
A raster plot shows individual events. A population rate smooths over many neurons or many trials to describe average activity. Rate models discard precise spike timing and keep a continuous activity variable. This is useful when the scientific question concerns stable states, oscillations, competition, or average input-output transformations of a circuit.

\paragraph{Formal definitions and notation.}
Let \(r_i(t)\geq0\) be the firing rate or activity of population \(i\). A standard rate model is
\[
\tau_i\frac{dr_i}{dt}=-r_i+\phi_i\!\left(\sum_j W_{ij}r_j+I_i(t)\right),
\]
where \(W_{ij}\) is the weight from population \(j\) to population \(i\), \(I_i(t)\) is external input, and \(\phi_i:\mathbb{R}\to\mathbb{R}_{\geq0}\) is a static transfer function. In vector form,
\[
\tau\dot r=-r+\phi(Wr+I),
\]
with componentwise nonlinearity \(\phi\).

For two excitatory-inhibitory populations, a Wilson-Cowan style model is
\[
\tau_E\dot r_E=-r_E+\phi_E(w_{EE}r_E-w_{EI}r_I+I_E),
\]
\[
\tau_I\dot r_I=-r_I+\phi_I(w_{IE}r_E-w_{II}r_I+I_I).
\]

\paragraph{Core mechanism: fixed points and gain.}
A fixed point \(r^*\) satisfies
\[
r^*=\phi(Wr^*+I).
\]
In one dimension,
\[
\tau\dot r=-r+\phi(wr+I).
\]
The fixed-point equation is \(r=\phi(wr+I)\). Linearizing around \(r^*\) gives
\[
\tau\frac{d}{dt}\delta r=\bigl(-1+w\phi'(wr^*+I)\bigr)\delta r.
\]
Thus the fixed point is locally stable if
\[
-1+w\phi'(wr^*+I)<0,
\]
or equivalently \(w\phi'(\cdot)<1\). Recurrent gain \(w\) and transfer slope \(\phi'\) together determine whether perturbations decay or grow.

\paragraph{Concrete example.}
Let \(\phi(u)=\max(0,u)\), \(w=0.5\), and \(I=4\). Since the fixed point is positive, \(r=\phi(0.5r+4)=0.5r+4\). Solving gives \(r^*=8\). The linear coefficient is \(-1+w\phi'= -1+0.5=-0.5\), so the fixed point is stable with relaxation time \(\tau/0.5=2\tau\).

If instead \(w=1.2\) and the ReLU remains active, the linear coefficient is \(0.2\), so perturbations grow in the one-dimensional model until some saturation not included in the ReLU approximation stops them. This shows why nonlinear saturation or inhibition is needed in strongly recurrent rate models.

\paragraph{Computational neuroscience example.}
Rate models describe cortical population dynamics such as excitation-inhibition balance, persistent activity, rivalry, and oscillations. A two-population E-I model can oscillate when excitation drives inhibition with a delay-like time constant, and inhibition then suppresses excitation. The variables are not individual membrane voltages; they are coarse activity levels of populations.

\paragraph{AI and machine-learning example.}
Continuous-time recurrent neural networks have the same form as rate models:
\[
\tau \dot h=-h+\sigma(Wh+Ux+b).
\]
The hidden state \(h\) is an artificial population activity. Stability analysis of \(Wh\) and the activation derivative predicts memory, saturation, exploding dynamics, or oscillations. This connects recurrent AI models to the dynamical-systems tools used for neural circuits.

\paragraph{Common misconceptions and failure modes.}
First, a rate is not a spike count; it has units of spikes per time or is a normalized activity proxy. Second, rate models can miss spike timing, synchrony, refractoriness, and correlations. Third, the transfer function \(\phi\) is not universal; it depends on the neuron model, input noise, and operating regime. Fourth, a stable rate fixed point does not imply every underlying spiking realization is identical or noise-free.

\paragraph{Exercises.}
\begin{enumerate}[label=\textbf{44.4.\arabic*.}, leftmargin=*]
\item \textbf{Mechanical.} In \(\tau\dot r=-r+\phi(wr+I)\), identify the leak term, recurrent input, external input, and transfer function.
\item \textbf{Proof or derivation.} Derive the one-dimensional local stability condition \(w\phi'(wr^*+I)<1\).
\item \textbf{Numerical/computational.} With \(\phi(u)=1/(1+e^{-u})\), \(w=2\), and \(I=0\), verify that \(r=0.5\) is not a fixed point. Then write the fixed-point equation that would need to be solved numerically.
\item \textbf{Interpretation.} What kinds of neural questions are better suited to a rate model than to a detailed spike-time model?
\item \textbf{Misconception/debugging.} A simulation of a rate model diverges to huge activity. Name two mathematical features that could prevent this in a more realistic model.
\end{enumerate}

\paragraph{Closing bridge.}
Point-neuron and rate models make spiking dynamics tractable at the scale of circuits. Chapter 45 adds the missing communication mechanism: synapses convert presynaptic spikes into postsynaptic currents, and plasticity rules change the strengths of those connections over time.

\clearpage
